\documentclass[11pt]{article}

\usepackage{amsmath,amssymb}
\usepackage{booktabs}
\usepackage{geometry}
\usepackage{listings}
\usepackage{url}

\lstset{frame=tb,
  aboveskip=3mm,
  belowskip=3mm,
  showstringspaces=false,
  columns=flexible,
  basicstyle={\small\ttfamily},
  numbers=none,
  breaklines=true,
  breakatwhitespace=true,
  tabsize=2
}

\geometry{margin=1in}

\title{An Optimized Time-Indexed ILP Formulation for Periodic Parallel Machine Scheduling with Resource Constraints}
\author{}
\date{}

\begin{document}

\maketitle

\section{System Model}\label{sec:model}

This note specifies the time-indexed integer linear program (ILP) built and
solved by \path{offline_phase/siren}. It is the model of
Section~VI of the paper (the ILP of SIREN), in the paper's notation, together
with the details the paper leaves to the implementation
(Section~\ref{sec:impl}). One model is built per candidate tuple, i.e., per
choice of one transmission configuration for every flow.

\subsection{Flows, Tasks and Resources}

The gNB serves $N$ periodic flows $\mathcal{F}=\{F_1,\dots,F_N\}$. A candidate
tuple fixes one configuration $K_i$ per flow (a segmentation option and a
combination of MCS, RBs and attempts, chosen in Section~V of the paper), so
that flow $F_i$ contributes $N_i$ identical TB-tasks
\[
\tau_{i,k}=(C_i, T_i, D_i, B_i), \qquad k=1,\dots,N_i,
\]
where:
\begin{itemize}
    \item $T_i$ is the period;
    \item $D_i$ is the relative deadline of the packet, with $D_i \le T_i$
          (constrained deadlines);
    \item $C_i = A_i\, t_\mathrm{tx}$ is the worst-case occupation time of one
          TB (its WCTT): the time it holds its HARQ process and its RBs,
          all $A_i$ attempts included;
    \item $B_i$ is the number of RBs the TB reserves at every TTI of $C_i$;
    \item $N_i$ is the number of TBs (segments) of one packet.
\end{itemize}
Releases are synchronous (the first packet of every flow is released at
time~$0$) and execution is non-preemptive. A packet is delivered
$t_\mathrm{up}$ after its last TB ends: the upper-layer delay holds neither a
HARQ process nor RBs, but it counts towards the packet's deadline.

Two resources are shared:
\begin{itemize}
    \item the set $\mathcal{M}$ of $M$ HARQ processes, identical processors:
          a TB occupies one of them, non-preemptively, for its whole $C_i$;
    \item a pool of $B^{\max}$ RBs, replenished at every TTI: the RBs of the
          TBs running in a TTI must total at most $B^{\max}$.
\end{itemize}

Time is slotted into TTIs of duration $\Delta$. The model works on TTI
indices: from here on, $T_i$, $D_i$, $C_i$, $t_\mathrm{up}$, $H$ and every
job time are integer numbers of TTIs, converted from the millisecond values
as described in Section~\ref{sec:impl} (exactly, for the paper's catalogue).
The hyperperiod and its TTIs are
\[
H = \operatorname{lcm}(T_1, \dots, T_N), \qquad
\mathcal{H} = \{0, 1, \dots, H-1\}.
\]

\subsection{Job Expansion}

Packet $r$ of flow $F_i$ is released at $r\,T_i$, for
$r = 0, \dots, H/T_i - 1$, and each of its $N_i$ TBs is a job. A job $j$ of
$F_i$ has
\[
r_j = r\,T_i, \qquad C_j = C_i, \qquad B_j = B_i, \qquad
d_j = r_j + D_i - t_\mathrm{up}.
\]
As in the paper, $d_j$ is the packet's absolute deadline less
$t_\mathrm{up}$: the latest time at which the job may \emph{end} for its packet
to be delivered on time. A job may therefore start at any
$t \in [r_j,\, d_j - C_j]$, its \emph{window}. The jobs of one packet share
$r_j$ and $d_j$, so the packet meets its deadline exactly when all its jobs
end by $d_j$.

After expansion, the periodic problem is a finite, discrete-time scheduling
problem over $[0, H)$ with the job set
\[
\mathcal{J} = \{J_1, \dots, J_Q\}, \qquad
Q = |\mathcal{J}| = \sum_{i=1}^{N} N_i \, \frac{H}{T_i}.
\]
We write $\mathcal{J}_{i,r} = \{j_1, \dots, j_{N_i}\}$ for the jobs of packet
$r$ of $F_i$, where $j_a$ is TB segment $a$.

\subsection{Notation Summary}

\subsubsection*{Sets}
\begin{center}
\begin{tabular}{ll}
\toprule
\textbf{Symbol} & \textbf{Description} \\
\midrule
$\mathcal{F}$ & Flows $F_1, \dots, F_N$ \\
$\mathcal{M}$ & HARQ processes, $|\mathcal{M}| = M$ \\
$\mathcal{J}$ & Jobs of one hyperperiod, $|\mathcal{J}| = Q$ \\
$\mathcal{J}_{i,r}$ & The $N_i$ jobs of packet $r$ of flow $F_i$ \\
$\mathcal{H}$ & TTI indices of one hyperperiod, $\{0, \dots, H-1\}$ \\
\bottomrule
\end{tabular}
\end{center}

\subsubsection*{Parameters (all times in TTIs)}
\begin{center}
\begin{tabular}{ll}
\toprule
\textbf{Symbol} & \textbf{Description} \\
\midrule
$T_i$, $D_i$ & Period and relative deadline of flow $F_i$ \\
$C_i$, $B_i$ & Occupation time (WCTT) and RBs per TTI of one TB of $F_i$ \\
$N_i$ & TBs per packet of $F_i$ \\
$t_\mathrm{up}$ & Upper-layer delay between a packet's last TB end and its delivery \\
$r_j$, $d_j$ & Release and latest end of job $j$, $d_j = r_j + D_i - t_\mathrm{up}$ \\
$C_j$, $B_j$ & Occupation time and RBs of job $j$ (those of its flow) \\
$B^{\max}$ & RBs available in every TTI \\
$M$ & Number of HARQ processes \\
$H$ & Hyperperiod, $\operatorname{lcm}(T_1, \dots, T_N)$ \\
$Q$ & Number of jobs, $\sum_i N_i H / T_i$ \\
$\Delta$ & TTI duration (milliseconds per index) \\
\bottomrule
\end{tabular}
\end{center}

\subsubsection*{Variables}
\begin{center}
\begin{tabular}{lll}
\toprule
\textbf{Symbol} & \textbf{Domain} & \textbf{Description} \\
\midrule
$x_{jpt}$ & $\{0,1\}$ & Job $j$ starts at TTI $t$ on HARQ process $p$ (Section~\ref{sec:vars}) \\
$u_p$ & $\{0,1\}$ & Process $p$ hosts a job (HARQ objective only, Section~\ref{sec:obj_harq}) \\
$B^{\mathrm{peak}}$ & $\{0, \dots, B^{\max}\}$ & Peak RB usage (RB-peak objective only, Section~\ref{sec:obj_rb}) \\
\bottomrule
\end{tabular}
\end{center}

\subsection{Remarks}
\begin{itemize}
    \item Releases are synchronous and execution is non-preemptive: once a job
          starts, it holds its HARQ process and its RBs for $C_j$ consecutive
          TTIs.
    \item Since $D_i \le T_i$, every job lies inside the hyperperiod:
          $d_j \le r_j + D_i \le r_j + T_i \le H$. No job crosses the boundary
          at $H$, so the table computed for $[0, H)$ can be replayed
          cyclically.
    \item The RB reservation is conservative: a job reserves $B_j$ RBs at every
          TTI of $C_j$, including TTIs in which it does not transmit
          (Section~IV-B of the paper). A table feasible under this
          reservation therefore also satisfies the real RB constraint.
\end{itemize}


%%%%%%%%%%%%%%%%%%%%%%%%%%%%%%%%%%%%%%%%%%%%%%%%%%%%%%%

\section{Decision Variables}\label{sec:vars}

The binary decision variable is
\[
x_{jpt}=
\begin{cases}
1, & \text{if job }j\text{ starts at TTI }t\text{ on HARQ process }p,\\
0, & \text{otherwise},
\end{cases}
\]
created for every $p \in \mathcal{M}$ and only for the start times of the
job's window,
\[
r_j \le t \le d_j - C_j .
\]
Variables outside the window are never created, which is what keeps the model
small (Section~\ref{sec:size}). A job that starts at $t$ occupies its process
and $B_j$ RBs during $[t, t + C_j)$, so it is running at TTI $t$ exactly when
its start lies in $[\mathcal{L}_j(t), \mathcal{R}_j(t)]$, with
\[
\mathcal{L}_j(t) = \max(r_j,\, t - C_j + 1), \qquad
\mathcal{R}_j(t) = \min(t,\, d_j - C_j).
\]
For notational convenience, the start time of job $j$ is
\[
S_j = \sum_{p\in\mathcal M} \sum_{t=r_j}^{d_j-C_j} t \, x_{jpt},
\]
which is well defined since each job starts exactly once
(constraint~\eqref{eq:job_once}).

%%%%%%%%%%%%%%%%%%%%%%%%%%%%%%%%%%%%%%%%%%%%%%%%%%%%%%%

\section{Feasibility Constraints}\label{sec:feas}

\subsection*{Assignment}

Each job starts exactly once, on one process, within its window:
\begin{equation}
\sum_{p\in\mathcal M} \sum_{t=r_j}^{d_j-C_j} x_{jpt} = 1,
\qquad \forall j\in\mathcal J.
\label{eq:job_once}
\end{equation}

\subsection*{HARQ Process Capacity}

Each HARQ process runs at most one job in every TTI:
\begin{equation}
\sum_{j\in\mathcal J} \sum_{t'=\mathcal{L}_j(t)}^{\mathcal{R}_j(t)} x_{jpt'} \le 1,
\qquad \forall p\in\mathcal M,\; t\in\mathcal H.
\label{eq:one_job_per_machine}
\end{equation}

\subsection*{RB Capacity}

The RBs of the jobs running in a TTI total at most $B^{\max}$:
\begin{equation}
\sum_{p\in\mathcal M} \sum_{j\in\mathcal J} B_j
\sum_{t'=\mathcal{L}_j(t)}^{\mathcal{R}_j(t)} x_{jpt'}
\le B^{\max},
\qquad \forall t\in\mathcal H.
\label{eq:max_resources}
\end{equation}

A sum in \eqref{eq:one_job_per_machine} or \eqref{eq:max_resources} is empty
when no job can be running at $t$; such a row is always satisfied and is not
generated.

%%%%%%%%%%%%%%%%%%%%%%%%%%%%%%%%%%%%%%%%%%%%%%%%%%%%%%%

\section{Symmetry-Breaking Constraints}\label{sec:sym}

The $N_i$ jobs of $\mathcal{J}_{i,r}$ have the same release, deadline,
occupation time and RBs, so any table stays valid when their starts and
processes are exchanged. To remove these equivalent solutions, consecutive
segments are ordered by start time:
\begin{equation}
S_{j_a} \le S_{j_{a+1}},
\qquad \forall\, \mathcal{J}_{i,r},\; a = 1, \dots, N_i - 1,
\label{eq:break_sym}
\end{equation}
or, in terms of the decision variables,
\[
\sum_{p\in\mathcal M} \sum_{t=r_{j_a}}^{d_{j_a}-C_{j_a}} t \, x_{j_a p t}
\le
\sum_{p\in\mathcal M} \sum_{t=r_{j_{a+1}}}^{d_{j_{a+1}}-C_{j_{a+1}}} t \, x_{j_{a+1} p t}.
\]
By transitivity, the $N_i - 1$ constraints per packet order all $N_i$
segments, instead of $\binom{N_i}{2}$ pairwise constraints. No schedule is
lost: relabelling the segments of any table in order of their start times
gives an equivalent table that satisfies \eqref{eq:break_sym}. Ties are
allowed, since two segments may start in the same TTI on different processes.

%%%%%%%%%%%%%%%%%%%%%%%%%%%%%%%%%%%%%%%%%%%%%%%%%%%%%%%

\section{Objective Functions}\label{sec:obj}

The formulation supports three objectives. Each keeps constraints
\eqref{eq:job_once} to \eqref{eq:break_sym}; only the last one changes the
right-hand side of \eqref{eq:max_resources}. Every objective value is an
integer, and the solver is run with a zero optimality gap
(Section~\ref{sec:impl}), so a reported optimum is exact.

\subsection{Feasibility (Default)}\label{sec:obj_feas}

Find any table that satisfies all constraints. The objective is a constant,
$\min 0$.

\subsection{Minimize HARQ Processes}\label{sec:obj_harq}

Introduce one binary variable per process,
\[
u_p =
\begin{cases}
1, & \text{if process }p\text{ hosts a job},\\
0, & \text{otherwise},
\end{cases}
\qquad p \in \mathcal{M},
\]
and solve
\begin{equation}
\min \sum_{p\in\mathcal M} u_p
\quad \text{s.t.} \quad
\sum_{j\in\mathcal J} \sum_{t=r_j}^{d_j-C_j} x_{jpt} \le Q \, u_p,
\qquad \forall p\in\mathcal M.
\label{eq:obj_harq}
\end{equation}
$Q = |\mathcal{J}|$ is a valid bound: at most $Q$ jobs start on one process.
The optimum is the number of HARQ processes the table uses.

\subsection{Minimize RB Peak}\label{sec:obj_rb}

Introduce one integer variable $B^{\mathrm{peak}} \in \{0, \dots, B^{\max}\}$
and replace the right-hand side of \eqref{eq:max_resources} with it:
\begin{equation}
\min B^{\mathrm{peak}}
\quad \text{s.t.} \quad
\sum_{p\in\mathcal M} \sum_{j\in\mathcal J} B_j
\sum_{t'=\mathcal{L}_j(t)}^{\mathcal{R}_j(t)} x_{jpt'}
\le B^{\mathrm{peak}},
\qquad \forall t\in\mathcal H.
\label{eq:obj_rb}
\end{equation}
Since $B^{\mathrm{peak}} \le B^{\max}$, \eqref{eq:obj_rb} implies
\eqref{eq:max_resources}, so the hard cap still holds. At an optimum,
$B^{\mathrm{peak}}$ equals the largest RB usage over the hyperperiod, and
$B^{\max} - B^{\mathrm{peak}}$ RBs are free in every TTI: the guaranteed
residual capacity for NRT traffic (RQ3 of the paper).

The paper states the same objective as \eqref{eq:max_resources} plus
$B^{\mathrm{peak}} \ge$ the RB usage of every TTI, without bounding
$B^{\mathrm{peak}}$. Both forms have the same optimal tables and the same
optimal value, because $B^{\mathrm{peak}} \le B^{\max}$ holds at any optimum
of the paper's form. $B^{\mathrm{peak}}$ is declared integer, so every
variable is integer and the model is a pure ILP. (It would be integral at the
optimum anyway, since every $B_j$ is an integer and every $x_{jpt}$ is
binary.)

%%%%%%%%%%%%%%%%%%%%%%%%%%%%%%%%%%%%%%%%%%%%%%%%%%%%%%%

\section{Model Size and Optimizations}\label{sec:size}

Let $w_i = D_i - t_\mathrm{up} - C_i + 1$ be the number of feasible starts of
a job of $F_i$ (its window width), and let
\[
\mathcal{H}^{c} = \bigcup_{j\in\mathcal J} [r_j, d_j) \subseteq \mathcal{H}
\]
be the TTIs in which some job can run. The feasibility model has exactly:

\begin{center}
\begin{tabular}{lll}
\toprule
\textbf{Item} & \textbf{Count} & \textbf{Source} \\
\midrule
Binary variables $x_{jpt}$ & $M \sum_{i=1}^{N} N_i \, \frac{H}{T_i} \, w_i$ & Section~\ref{sec:vars} \\
Assignment rows & $Q$ & \eqref{eq:job_once} \\
HARQ-capacity rows & $M \, |\mathcal{H}^{c}|$ & \eqref{eq:one_job_per_machine} \\
RB-capacity rows & $|\mathcal{H}^{c}|$ & \eqref{eq:max_resources} \\
Symmetry-breaking rows & $\sum_{i=1}^{N} \frac{H}{T_i} \, (N_i - 1)$ & \eqref{eq:break_sym} \\
\bottomrule
\end{tabular}
\end{center}

Without the window restriction there would be $M Q H$ binary variables; with
it, each job has $M w_i$, independent of $H$. Symmetry breaking adds rows but
no variables. Every count grows linearly with the number of jobs, hence with
$H$ for a fixed set of periods. The objectives add:

\begin{itemize}
    \item \textbf{Feasibility}: nothing;
    \item \textbf{HARQ processes}: $M$ binary variables $u_p$ and $M$ rows
          \eqref{eq:obj_harq};
    \item \textbf{RB peak}: one integer variable $B^{\mathrm{peak}}$ and no
          row (each RB-capacity row gains one coefficient).
\end{itemize}

The \texttt{num\_variables} and \texttt{num\_constraints} fields of the report
(Section~\ref{sec:impl}) count exactly these variables and rows.

%%%%%%%%%%%%%%%%%%%%%%%%%%%%%%%%%%%%%%%%%%%%%%%%%%%%%%%

\section{Implementation Notes}\label{sec:impl}

\subsection*{Input}

The input is the planner's \texttt{planner\_report.json}: per flow, the
configurations of Section~V of the paper, grouped by segmentation option.
An excerpt (one flow, one segmentation option with two configurations):

\begin{lstlisting}
{
  "parameters": {
    "a_max": 4, "rlc_header_bytes": 8, "t_tx_ms": 1.0, "t_up_ms": 4.0
  },
  "applications": [
    {
      "app": "DiscreteAutomation_T30",
      "period": 30.0,
      "deadline": 15.0,
      "segmentation_options": [
        {
          "num_segments": 2,
          "tbs": 469,
          "mbc_combinations": [
            { "a": 2, "b": 64, "bler": 0.006375, "c": 2.0, "m": 2 },
            { "a": 3, "b": 50, "bler": 0.024140625, "c": 3.0, "m": 3 }
          ]
        }
      ]
    }
  ]
}
\end{lstlisting}

The model reads the following fields; times are in milliseconds and are
converted to TTIs of $\Delta$~milliseconds (option \texttt{--tti}):

\begin{center}
\begin{tabular}{lll}
\toprule
\textbf{Field} & \textbf{Symbol} & \textbf{Conversion} \\
\midrule
\texttt{parameters.t\_up\_ms} & $t_\mathrm{up}$ & $\lceil t_\mathrm{up} / \Delta \rceil$; required, the run is refused without it \\
\texttt{period} & $T_i$ & $T_i / \Delta$, which must be an integer \\
\texttt{deadline} & $D_i$ & $\lfloor D_i / \Delta \rfloor$ \\
\texttt{num\_segments} & $N_i$ & none \\
\texttt{mbc\_combinations[].b} & $B_i$ & none \\
\texttt{mbc\_combinations[].c} & $C_i$ & $\lceil C_i / \Delta \rceil$ ($C_i = A_i\, t_\mathrm{tx}$) \\
\bottomrule
\end{tabular}
\end{center}

Each rounding goes the way that cannot make an infeasible tuple look feasible:
a deadline loses slack, an occupation time and the upper-layer delay grow. A
period off the grid aborts the run (exit status 1), since the hyperperiod
would not be a feasibility interval. The fields \texttt{app}, \texttt{m},
\texttt{a} and \texttt{tbs} only label the output (\texttt{app} and
\texttt{m} are nevertheless required). The other fields are ignored.

\subsection*{Candidate Tuples}

The configurations of flow $F_i$ (every \texttt{mbc\_combinations} entry of
every segmentation option) form its set $\mathcal{P}_i$. A candidate tuple
picks one configuration per flow, and each of the $\prod_i |\mathcal{P}_i|$
tuples is a separate model. The tuples are enumerated like an odometer (the
last flow changes fastest) and numbered from 1; this number is the
\texttt{combo\_index} of the output. The sweep does not stop at the first
schedulable tuple: it decides every tuple, or as many as the time budget
allows.

If some flow has no configuration ($\mathcal{P}_i = \emptyset$), no model is
built: the report lists the flow under \texttt{infeasible\_flows}, and the run
still exits with status 0, since this is a verdict on the flowset and not a
failure.

Two checks reject a tuple without building a model, and report it as not
schedulable: some $D_i > T_i$ (outside the constrained-deadline model of
Section~\ref{sec:model}), or some empty window, $C_i + t_\mathrm{up} > D_i$.

\subsection*{Solver}

The model is passed to Gurobi in two calls (\texttt{GRBnewmodel} for the
variables, \texttt{GRBXaddconstrs} for every constraint), in one Gurobi
environment started once per run, with:
\begin{itemize}
    \item one thread (\texttt{Threads=1}): the experiments run one process per
          core, and the solvers' own threads would otherwise starve each
          other;
    \item \texttt{MIPGap = MIPGapAbs = 0}: every objective is an integer, so
          the optimum is closed exactly;
    \item \texttt{TimeLimit} set to the time this tuple may use (see
          Timeouts below), at least 1~s.
\end{itemize}
Only two outcomes decide a tuple:
\begin{center}
\begin{tabular}{p{0.3\linewidth}p{0.62\linewidth}}
\toprule
\textbf{Gurobi status} & \textbf{Verdict} \\
\midrule
\texttt{OPTIMAL} & schedulable: the table is read from the $x_{jpt} > 0.5$ \\
\texttt{INFEASIBLE}, \texttt{INF\_OR\_UNBD} & not schedulable (every variable is bounded, so both mean infeasible) \\
\texttt{TIME\_LIMIT}, \texttt{ITERATION\_LIMIT} & timed out: undecided, even if a feasible solution was found \\
anything else, or an API error & solver error: undecided \\
\bottomrule
\end{tabular}
\end{center}
An undecided tuple never makes a flowset proved unschedulable.

\subsection*{Timeouts}

Two limits apply, both wall-clock:
\begin{itemize}
    \item \texttt{--timeout} $S$: the budget of the whole sweep of one flowset
          (default 300~s; 0 means no budget). When it runs out, the sweep
          stops and the report says \texttt{"timed\_out": true}.
    \item \texttt{--candidate-timeout} $c$: the cap of one tuple (default 0,
          i.e., off). A tuple that reaches it is listed in
          \texttt{timed\_out\_combos} and the sweep moves on to the next one.
\end{itemize}
Each solve gets the smaller of the remaining budget and $c$. The paper's
experiment uses no sweep budget and $c = 60$~s
(\texttt{scripts/sweep\_experiment.sh}).

\subsection*{Objective Selection}

\begin{itemize}
    \item \texttt{--objective feasibility}: Section~\ref{sec:obj_feas}
          (default);
    \item \texttt{--objective harq}: Section~\ref{sec:obj_harq}
          (\texttt{machines} is accepted as an alias);
    \item \texttt{--objective rb\_peak}: Section~\ref{sec:obj_rb}, searched
          over the tuples according to \texttt{--rb-peak-search}:
    \begin{itemize}
        \item \texttt{full} (the default of the binary): every tuple is
              solved with the RB-peak objective;
        \item \texttt{bounded} (the default of the experiment scripts):
              every tuple is first decided with the feasibility objective.
              The schedulable ones are then re-solved with the RB-peak
              objective in ascending order of their mean RB usage
              $\bar{u} = \sum_i N_i B_i C_i / T_i$. Since no table's peak is
              below its mean usage, the search stops once the best peak found
              is at most $\lceil \bar{u} \rceil$ of the next tuple. A tuple
              that times out or fails in this second pass keeps its
              feasibility table, with no peak; the best peak is optimal over
              the decided tuples only when no such tuple occurred and the
              budget held (\texttt{optimal\_over\_decided}).
    \end{itemize}
\end{itemize}
The paper's experiment uses \texttt{rb\_peak} with the \texttt{bounded}
search.

\subsection*{Output}

The report is written to the path given by \texttt{--output-scheduling}, or
to standard output. It holds the configuration of the run, the number of
tuples (total, tested, schedulable, timed out, errored, with the indices of
the last two), \texttt{infeasible\_flows}, one entry per schedulable tuple
(its flows, hyperperiod, \texttt{num\_jobs}, \texttt{num\_variables},
\texttt{num\_constraints}, solve time, and \texttt{harq\_used} or
\texttt{rb\_peak\_used}), and, for the bounded search, a summary of the second
pass.

Every schedulable tuple also gets a dispatch table,
\texttt{combo\_<index>.json} in \texttt{--schedule-dir} (default
\texttt{schedules}; \texttt{--no-schedules} skips them). It lists every job
with its flow, segment, HARQ process, release, start, end and absolute
deadline, in TTIs and in milliseconds, and its RBs. The end is when the TB
frees its process and RBs; the packet is delivered $t_\mathrm{up}$ after the
end of its last TB, and that is what the deadline bounds. With the bounded
search, the RB-peak table replaces the feasibility table of the same tuple.
Every field is described in \texttt{doc/output\_formats.md}.

\subsection*{Command Line}

\begin{center}
\begin{tabular}{lll}
\toprule
\textbf{Option} & \textbf{Meaning} & \textbf{Default} \\
\midrule
\texttt{--n-rb} & $B^{\max}$ & 50 \\
\texttt{--m-harq} & $M$ & 2 \\
\texttt{--tti} & $\Delta$, in milliseconds & 0.25 \\
\texttt{--objective} & \texttt{feasibility}, \texttt{harq} or \texttt{rb\_peak} & \texttt{feasibility} \\
\texttt{--rb-peak-search} & \texttt{full} or \texttt{bounded} & \texttt{full} \\
\texttt{--timeout} & sweep budget, seconds (0: none) & 300 \\
\texttt{--candidate-timeout} & per-tuple cap, seconds (0: off) & 0 \\
\texttt{--output-scheduling} & report path & standard output \\
\texttt{--schedule-dir} & dispatch-table directory & \texttt{schedules} \\
\texttt{--no-schedules} & write no dispatch table & off \\
\texttt{--quiet}, \texttt{-q} & no per-tuple progress lines & off \\
\bottomrule
\end{tabular}
\end{center}

\end{document}
